\section{Scientific conclusions suitable for an article}
\label{sec:report-article-conclusions}
The completed analysis establishes a reproducible mathematical diagnosis of the frozen reconstruction. It does not establish a new validated biological model. The most defensible central conclusion is that explicit conservation and acceptable reference physiology do not suffice to reproduce the AE4 secretion phenotype: the adopted network largely replaces the missing AE4 chloride input, while the prescribed apical recruitment construction recovers the selected magnitude only under an additional strong and unvalidated coupling.

This conclusion is supported by several complementary results. The conserved amount formulation identifies the alkalinity cost of productive AE4 transport. The exact chloride identity identifies which explicit source contributions survive in a stationary balance. The full equilibrium and transient calculations quantify the state mediated response that the identity alone cannot determine. The parameter audit exposes which effective settings are independently supported and which remain assumptions. Finally, the inverse calculation separates mathematical sufficiency of a chosen coupling from biological evidence for that coupling. None of those steps can be replaced by a single comparison with the experimental secretion percentage.

\subsection{What the model explains internally}
Within the frozen configuration, a fall in the AE4 chloride contribution is substantially offset by the rest of the transport network, particularly NKCC1. The local equilibrium gain, the integrated replacement fraction and the relative increase in NKCC1 loading all support this interpretation, provided their distinct definitions are retained. The small stationary effect and the larger, but still small, 600 s cumulative effect are compatible with the presence of changing ionic stores and a finite lumen. They are not competing estimates of one observable.

The model also explains why alkalinity must be examined together with carbon and volume. A sustained AE4 demand must be supplied by the admissible alkalinity sources or by depletion of a finite store. An NBC absent root with much smaller AE4 loading is consistent with that requirement. It does not invalidate the budget. Conversely, observing NBC support in one WT trajectory does not establish a universal necessity theorem or identify the molecular pathway.

The local mode analysis shows that the system already possesses dynamics on a scale of several minutes. Its modal contributions to secretion are not all zero, and their signs can cancel in the stationary sensitivity. Accordingly, a failed temporal phenotype should not be attributed to absence of a slow state without analysing how the perturbation excites the existing modes and how secretion observes them. The documented correction of the regulatory clamp derivative is a numerical differentiation correction, not a change to the production mechanism.

\subsection{What fails as an account of the biological phenotype}
The reference finite protocol does not reproduce the large sustained secretion loss described in the source experimental comparison. The compensation responsible for much of this attenuation has not been independently established in the corresponding physiological preparation. In addition, the substitution of source vectors into inherited scalar AE4 kinetics leaves explicit thermodynamic qualifications. These are substantive limitations, not typesetting or numerical convergence problems.

The Task 41 family illustrates a distinct limitation. It can approach the selected finite duration deficit while retaining the declared 600 s gates, but it requires most fully stimulated CaCC recruitment to depend on AE4 in that particular formula. The original temporal and ionic discrepancies remain, and the separate extended null trajectories later cross the upper pH gate. Thus a finite duration fit to one chosen magnitude does not establish a sustained, independently validated secretion mechanism.

The source analysis does not isolate NKCC1 as the uniquely incorrect component. The existing rational law already has finite forward and reverse limits. Its strong compensatory response could depend on its effective capacity or stimulation, on ionic state dependence, on the pump and bicarbonate pathways, on the common initial condition, or on missing coupling elsewhere. Establishing the responsible defect requires a new comparison of mechanistically distinct candidates and matched experimental protocols. The results here justify that comparison and identify quantities to measure; they do not predetermine its winner.
\sourcefile{evidence/task44/integrated_context.tex; evidence/task44/literature_comparison.tex; evidence/task44/inverse_results.tex; evidence/task43/report.tex}

\section{Claims that must remain qualified}
\label{sec:report-qualified-claims}
\begin{longtable}{>{\raggedright\arraybackslash}p{.30\linewidth}>{\raggedright\arraybackslash}p{.63\linewidth}}
\toprule
Claim to avoid & Supported formulation\\\midrule\endhead
The full model is biologically correct because all checks pass. & Conservation, numerical agreement, physiological gates and thermodynamic compatibility are separate tests. Some physiological mathematical roots retain an opposed source cycle. Numerical replay does not validate biology.\\
NKCC1 is proved to be the sole defective mechanism. & The frozen network compensates strongly, and several poorly constrained settings affect that response. Attribution to a unique defective component remains unresolved.\\
NKCC1 must never increase after AE4 loss. & The source reports no detectable increase in the particular isolated initial uptake assay. That observation must be compared with an appropriate assay model; it is not an identity across all states and protocols.\\
The chloride cancellation proves no AE4 effect. & It removes an explicit source term from one stationary balance. The state mediated derivative remains and is quantified locally.\\
The rank result proves whole cell observations cannot distinguish source classes. & The five specified joint source projections have rank five. Whether experimental observables identify a source vector also depends on kinetics, excitation, measurement and uncertainty.\\
NBC is necessary for all physiological secretion. & The obstruction is conditional on a specified AE4 loading demand and admissible support bounds. The source includes an NBC absent physiological root at substantially smaller AE4 loading.\\
A missing slow regulator is proved necessary. & Slow local modes already occur and project onto secretion. Their existence alone does not reproduce the experimental temporal signature or exclude missing regulation.\\
The inverse crossing is a global minimum or a confidence bound. & It is a first numerical crossing in one ordered and locally refined family. The 30.3\% comparison equals a reported mean minus one standard error.\\
A local sensitivity proves robustness over physiological variation. & Derivatives are verified near a reference point. The audit supplies no defensible uncertainty intervals for many effective parameters, so finite range robustness is not established.\\
The 5\% model case independently validates an expression experiment. & It is a prescribed expression perturbation in this source record. Independent experimental status requires an appropriate source, not the presence of a simulated curve.\\
All failed historical numerical attempts have been reconstructed. & Four early unsuccessful continuation attempts survive only as an aggregate statement. Individual seeds and messages were unavailable and have not been invented.\\\bottomrule
\end{longtable}

\section{What must be constrained in a subsequent reconstruction}
\label{sec:report-future}
The source sensitivities prioritise functional pump capacity and apical distribution, effective NKCC1 capacity and stimulation, NBC capacity and saturation, NHE1 abundance and recruitment, the finite buffer pool and CaCC conductance or recruitment. These priorities are local to the analysed configuration. Their importance does not supply their physiological ranges, and a small sensitivity does not prove that a parameter is unimportant in all regimes.

The most informative measurements would separate transporter abundance from its state dependence and measure relevant quantities under comparable stimulation and ionic conditions. The source report identifies combined sodium and bicarbonate flux, current and pH measurements for alkalinity support; ion conditioned NKCC, NHE and pump measurements; simultaneous total inorganic carbon, buffering, chloride and volume time courses; and matched stimulated CaCC current or surface recruitment during acute AE4 perturbation. Combined cation, chloride, carbon and reversal measurements are needed to discriminate proposed AE4 source cycles and compatible kinetic implementations. These are proposed discriminating experiments, not data already available in this analysis.

The comparison of acute perturbations with established knockout preparations also requires an explicit initial state model. Reusing the inherited WT resting state across expression cases is a controlled computational intervention, but it does not reproduce every resting adaptation present in a knockout gland. Likewise, a single acinar cell with a finite lumen does not include ductal modification or an entire gland. Those differences should be modelled or bounded when they matter to an experimental comparison, not hidden by equating all secretion observables.

A subsequent model should therefore explain both the observed output and the mechanism that changes it, while preserving the conserved balances. It should distinguish data used for construction and selection from observations left for independent testing. If an observation has already been used to guide candidate selection, it is no longer an untouched validation test. This is a methodological consequence of the diagnosis, rather than an additional result proved by Tasks 43 and 44.

\section{Completeness, source limits and reproducible use}
\label{sec:report-completeness}
The report contains the complete final Task 44 technical exposition and the complete Task 43 provenance report, with explicit source namespaces. It expands the numerical tables beyond those printed in the original 26 and 11 page reports. All 26 state vectors, all 26 dynamic Jacobian matrices and all eigenvalues are printed. Every inventory record is given its own detailed entry, and all extended parameter metrics and both inverse threshold sensitivity observables are tabulated.

The electronic evidence is part of the report, not an optional source of convenient extra facts. It preserves full precision numbers, complete modal vectors, every stored nearby trial, source locations, diagnostic residuals, failure records, verification exclusions and original scripts. The PDF rounds table values for reading, and the source manifest supplies hashes of the unchanged underlying files. A reader can follow the report to a particular source record and then to the exact pinned commit. The build does not depend on live branch heads or regenerate scientific simulations.

Some detail is unavailable in the source record, most notably the individual seeds and errors of the four historical continuation failures and experimentally justified ranges for many effective parameters. No report can legitimately fill those gaps by inventing plausible values. They are therefore stated explicitly. The source literature review also distinguishes full article inspection from abstract or caption level access. This compilation does not silently upgrade that review to a systematic search of every later experiment.

The final technical package is a basis for rigorous article development, not a declaration that all scientific questions are closed. It preserves the results that survive, the conditional assumptions that make them true, the calculations that contradict earlier interpretations, and the evidence needed to test a better model. That is the contribution of the completed mathematical analysis.
\sourcefile{evidence/task44/FINAL_STATUS.md; evidence/task43/FINAL_STATUS.md; SOURCE_MANIFEST.json; COVERAGE.md}
